% GENERATED FILE. Edit canonical lessons, then run npm run generate:books.
% canonical-source-hash: fnv1a64:a40273e389cdad30
% canonical-lessons: RU-C200-sekret, RU-C200-shutka, RU-C200-smekh, RU-C200-slyozy, RU-C200-ulybka

\chapter{Secret, Joke, Laughter, Tears, Smile}
\label{ch:ru-secret-joke-laughter-tears-smile}
\hlchaptermodality{200}

\begin{chapteropening}
\textbf{By the end of this chapter:} \emph{I can say a secret, a joke, laughter, tears and a smile.}

Complete the last lesson of chapter 200: I can say a secret, a joke, laughter, tears and a smile.
\end{chapteropening}

\section[\texorpdfstring{\ru{секрет} (sekrét)}{секрет (sekrét)}]{\ru{секрет} (sekrét) — a secret}

\label{lesson:RU-C200-sekret}



\begin{quote}
Before the new one: say the Russian for a tourist, then the Russian for a text.
\end{quote}

\subsection*{You'll want to know: \ru{секрет}}
\textbf{\ru{секрет}} (\emph{sekrét}) — \textquotedblleft{}a secret\textquotedblright{}.

\begin{cousinweb}
\textbf{\ru{секрет}} — from Peter the Great's time; it perhaps came through Polish \emph{sekret}, or from French \emph{secret}, from Latin \emph{sēcrētus}, \textquotedblleft{}set apart, hidden\textquotedblright{} (Vasmer).
\end{cousinweb}

\begin{grammarlens}[title={What you've built}]
One more word for this chapter.
\end{grammarlens}

\subsection*{Your turn}
\begin{itemize}
\raggedright
  \item \emph{Say it:} \emph{sekrét}
  \item \emph{Say it:} \emph{sekrét}, once more
  \item \emph{Say it:} say \emph{sekrét}
  \item \emph{Recall:} say the Russian for a tourist, then the Russian for a text, then say \emph{sekrét} again
\end{itemize}

\begin{tcolorbox}[breakable,colback=teal!4,colframe=teal!35!black,title={Before you move on}]
What does \ru{секрет} mean? (\textquotedblleft{}A secret\textquotedblright{}.) Say it once more.
\end{tcolorbox}
\section[\texorpdfstring{\ru{шутка} (shútka)}{шутка (shútka)}]{\ru{шутка} (shútka) — a joke}

\label{lesson:RU-C200-shutka}



\begin{quote}
Before the new one: say the Russian for a text, then the Russian for a secret.
\end{quote}

\subsection*{You'll want to know: \ru{шутка}}
\textbf{\ru{шутка}} (\emph{shútka}) — \textquotedblleft{}a joke\textquotedblright{}.

\begin{cousinweb}
\textbf{\ru{шутка}} comes from \emph{\ru{шутить}}, \textquotedblleft{}to joke\textquotedblright{}, and \emph{\ru{шут}}, \textquotedblleft{}a jester\textquotedblright{}, recorded from 1076. Vasmer reports a suggested link with Lithuanian \emph{siaũsti}, \textquotedblleft{}to rage, to frolic\textquotedblright{}, and doubts a loan from German \emph{Schaute}, \textquotedblleft{}a fool\textquotedblright{}.
\end{cousinweb}

\begin{grammarlens}[title={What you've built}]
One more word for this chapter.
\end{grammarlens}

\subsection*{Your turn}
\begin{itemize}
\raggedright
  \item \emph{Say it:} \emph{shútka}
  \item \emph{Say it:} \emph{shútka}, once more
  \item \emph{Say it:} say \emph{shútka}
  \item \emph{Recall:} say the Russian for a text, then the Russian for a secret, then say \emph{shútka} again
\end{itemize}

\begin{tcolorbox}[breakable,colback=teal!4,colframe=teal!35!black,title={Before you move on}]
What does \ru{шутка} mean? (\textquotedblleft{}A joke\textquotedblright{}.) Say it once more.
\end{tcolorbox}
\section[\texorpdfstring{\ru{смех} (smekh)}{смех (smekh)}]{\ru{смех} (smekh) — laughter}

\label{lesson:RU-C200-smekh}



\begin{quote}
Before the new one: say the Russian for a secret, then the Russian for a joke.
\end{quote}

\subsection*{You'll want to know: \ru{смех}}
\textbf{\ru{смех}} (\emph{smekh}) — \textquotedblleft{}laughter\textquotedblright{}.

\begin{cousinweb}
\textbf{\ru{смех}} $\leftarrow$ Proto-Slavic \emph{*směxŭ}, with \emph{\ru{смеяться}}, \textquotedblleft{}to laugh\textquotedblright{} — inherited (Vasmer). Vasmer relates the root to Sanskrit \emph{smáyate}, \textquotedblleft{}smiles\textquotedblright{}, and Middle English \emph{smilen}, \textquotedblleft{}to smile\textquotedblright{}.
\end{cousinweb}

\begin{grammarlens}[title={What you've built}]
One more word for this chapter.
\end{grammarlens}

\subsection*{Your turn}
\begin{itemize}
\raggedright
  \item \emph{Say it:} \emph{smekh}
  \item \emph{Say it:} \emph{smekh}, once more
  \item \emph{Say it:} say \emph{smekh}
  \item \emph{Recall:} say the Russian for a secret, then the Russian for a joke, then say \emph{smekh} again
\end{itemize}

\begin{tcolorbox}[breakable,colback=teal!4,colframe=teal!35!black,title={Before you move on}]
What does \ru{смех} mean? (\textquotedblleft{}Laughter\textquotedblright{}.) Say it once more.
\end{tcolorbox}
\section[\texorpdfstring{\ru{слёзы} (slyózy)}{слёзы (slyózy)}]{\ru{слёзы} (slyózy) — tears}

\label{lesson:RU-C200-slyozy}



\begin{quote}
Before the new one: say the Russian for a joke, then the Russian for laughter.
\end{quote}

\subsection*{You'll want to know: \ru{слёзы}}
\textbf{\ru{слёзы}} (\emph{slyózy}) — \textquotedblleft{}tears\textquotedblright{}.

\begin{cousinweb}
\textbf{\ru{слёзы}} is the plural of \emph{\ru{слеза}}, \textquotedblleft{}a tear\textquotedblright{}, Proto-Slavic \emph{*slĭza} — inherited. Vasmer thinks it more likely belongs with \emph{\ru{слизь}}, \textquotedblleft{}slime\textquotedblright{}, and Middle Low German \emph{slîk}, \textquotedblleft{}slime\textquotedblright{}.
\end{cousinweb}

\begin{grammarlens}[title={What you've built}]
One more word for this chapter.
\end{grammarlens}

\subsection*{Your turn}
\begin{itemize}
\raggedright
  \item \emph{Say it:} \emph{slyózy}
  \item \emph{Say it:} \emph{slyózy}, once more
  \item \emph{Say it:} say \emph{slyózy}
  \item \emph{Recall:} say the Russian for a joke, then the Russian for laughter, then say \emph{slyózy} again
\end{itemize}

\begin{tcolorbox}[breakable,colback=teal!4,colframe=teal!35!black,title={Before you move on}]
What does \ru{слёзы} mean? (\textquotedblleft{}Tears\textquotedblright{}.) Say it once more.
\end{tcolorbox}
\section[\texorpdfstring{\ru{улыбка} (ulýbka)}{улыбка (ulýbka)}]{\ru{улыбка} (ulýbka) — a smile}

\label{lesson:RU-C200-ulybka}



\begin{quote}
Before the new one: say the Russian for laughter, then the Russian for tears.
\end{quote}

\subsection*{You'll want to know: \ru{улыбка}}
\textbf{\ru{улыбка}} (\emph{ulýbka}) — \textquotedblleft{}a smile\textquotedblright{}.

\begin{cousinweb}
\textbf{\ru{улыбка}} goes with \emph{\ru{улыбаться}}, \textquotedblleft{}to smile\textquotedblright{}, from dialect \emph{\ru{лыбить}}, \textquotedblleft{}to smile\textquotedblright{}. Vasmer calls \emph{\ru{лыбить}} unclear and finds the proposed links, with \emph{\ru{луб}}, \textquotedblleft{}bast\textquotedblright{}, or Latin \emph{labō}, \textquotedblleft{}I totter\textquotedblright{}, unconvincing.
\end{cousinweb}

\begin{grammarlens}[title={What you've built}]
One more word for this chapter.
\end{grammarlens}

\subsection*{Your turn}
\begin{itemize}
\raggedright
  \item \emph{Say it:} \emph{ulýbka}
  \item \emph{Say it:} \emph{ulýbka}, once more
  \item \emph{Say it:} say \emph{ulýbka}
  \item \emph{Recall:} say the Russian for laughter, then the Russian for tears, then say \emph{ulýbka} again
\end{itemize}

\begin{tcolorbox}[breakable,colback=teal!4,colframe=teal!35!black,title={Before you move on}]
What does \ru{улыбка} mean? (\textquotedblleft{}A smile\textquotedblright{}.) Say it once more.
\end{tcolorbox}
